% Canonical Paper IV section.
The preceding examples varied semantic quotients and cut order.  We now keep
the final map fixed and vary its realization.  Matrix parenthesization changes
arithmetic work and intermediate shape; reverse-mode checkpointing exchanges
stored history for recomputation.

\subsection{Matrix-chain parenthesization}

Let
\[
  A\in K^{d_0\times d_1},
  \qquad B\in K^{d_1\times d_2},
  \qquad C\in K^{d_2\times d_3}.
\]
Both $(AB)C$ and $A(BC)$ evaluate to the same linear map whenever the products
are defined.

\begin{proposition}[Three-matrix work and principal intermediates]
\label{prop:matrix-chain-costs}
Under ordinary dense multiplication, counting scalar multiplications,
\[
  W_L=d_0d_1d_2+d_0d_2d_3,
  \qquad
  R_L=d_0d_2,
\]
\[
  W_R=d_1d_2d_3+d_0d_1d_3,
  \qquad
  R_R=d_1d_3,
\]
where $R_L$ and $R_R$ are the numbers of scalar entries in the principal
intermediate matrices $AB$ and $BC$.
\end{proposition}

\begin{proof}
Multiplying an $r\times s$ matrix by an $s\times t$ matrix uses $rst$ scalar
multiplications in the stated algorithm and produces an $r\times t$ matrix.
Apply this twice in each order.
\end{proof}

For $(d_0,d_1,d_2,d_3)=(100,10,1,2)$,
\[
  (W_L,R_L)=(1200,100),
  \qquad
  (W_R,R_R)=(2020,20).
\]
The lower-work plan has the larger principal intermediate.  The $R$ values are
not complete peak-memory formulas: input retention, output allocation,
overwrite policy, blocking, and temporaries still need to be declared.

For four matrices, the five parenthesizations form the vertices of an
associahedral pentagon.  If an edge is labelled by the difference of endpoint
costs, Proposition~\ref{prop:additive-exactness-obstruction} makes its sum around
the pentagon zero.  Positive work spent traversing the pentagon is action, not
holonomy.

\subsection{A chain and its reverse sweep}

Consider a length-$N$ composition
\[
  x_{i+1}=f_i(x_i),
  \qquad 0\le i<N,
\]
and its reverse accumulation
\[
  \bar x_i=Df_i(x_i)^*\bar x_{i+1}.
\]
The reverse step requires the corresponding forward state $x_i$.

Assume each forward evaluation and each reverse local step has unit cost and
each state has fixed size.

\begin{proposition}[Two exact endpoint schedules]
\label{prop:reverse-chain-extremes}
The following schedules compute the same reverse result.
\begin{enumerate}
  \item Store all forward states.  The total local work is $2N$ and the number
  of stored forward states is $\Theta(N)$.
  \item Store only the input and recompute each required prefix.  The total
  local work is
  \begin{equation}
    2N+\frac{N(N-1)}2,
    \label{eq:naive-recompute-work}
  \end{equation}
  and only $\Theta(1)$ states need be live, apart from the declared input,
  output, and control records.
\end{enumerate}
\end{proposition}

\begin{proof}
The store-all schedule performs $N$ forward and $N$ reverse local operations.
For the second schedule, the initial forward pass still costs $N$, the reverse
steps cost $N$, and reconstructing $x_i$ from $x_0$ costs $i$ forward steps.
Summing $i=0,\ldots,N-1$ gives Equation~\eqref{eq:naive-recompute-work}.
\end{proof}

Checkpointing interpolates between these extremes.  Under a specified number
of checkpoints and a uniform linear chain, schedules such as
\textsc{Revolve} optimize the recomputation pattern for their declared model
and objective~\cite{GriewankWalther2000}.

\subsection{History condensation and re-expansion}

A checkpoint is a materialized summary of the forward history sufficient to
restart from that point.  Erasing it condenses the live trace further;
recomputation re-expands the missing segment from an earlier checkpoint.  This
interpretation is structurally compatible with AEG's process-before-result
principle, but it does not identify a nonlinear derivative pullback with one
global projective inverse.

The example validates the live-configuration correction of
Theorem~\ref{thm:completed-set-insufficient}.  Both schedules eventually compute
the same forward nodes and the same reverse result, yet their live states and
work differ.  No monotone set of ever-computed nodes can encode that tradeoff.

\subsection{Reversible computation boundary}

A logically reversible realization must retain or reconstruct information that
a many-to-one step would otherwise erase.  General reversible simulations
therefore exhibit time--space tradeoffs; Bennett's pebbling construction is a
classical model~\cite{Bennett1989}.  The present paper uses this as an external
calibration, not as a theorem that every AEG condensation has the same optimal
tradeoff.

\subsection{What equal semantics permits us to conclude}

Matrix chains and checkpoint schedules prove that one semantic map can have
incomparable resource vectors.  They support the Pareto viewpoint of
Equation~\eqref{eq:pareto-resource-geometry}.  They do not show that all time
and space costs are coordinates of one scalar ``representation volume.''  A
unification would require a model-specific theorem relating semantic residuals,
live encodings, and allowed re-expansions.
